% !PRINTED 336
\weberpage{336}
%<<CONT
und höhere Potenzen von $h$ vorkommen, einen kleinen Werth haben. Es
wird also unter den geeigneten Voraussetzungen
\weq{4}{ \alpha' = \alpha - \frac{f(\alpha)}{f'(\alpha)} }
als eine bessere Annäherung an den wahren Werth der Wurzel zu betrachten
sein.

Ersetzt man dann $\alpha$ durch $\alpha'$, so wird man in
\[ \alpha'' = \alpha' - \frac{f(\alpha')}{f'(\alpha')} \]
eine noch bessere Näherung erhalten u.s.f.

Es bleiben hier aber noch folgende beiden Fragen zu beantworten:
\begin{enumerate}
\item[1.] \gesperrt{Unter welchen Voraussetzungen ist $\alpha'$
  wirklich ein besserer Werth als $\alpha$?}
\item[2.] \gesperrt{Wie kann man den Grad der Genauigkeit schätzen, den
  man so erreicht?}
\end{enumerate}
Diese Fragen hat \gesperrt{Fourier} beantwortet; er macht aber dabei
folgende Voraussetzung:

\gesperrt{Es ist eine Wurzel von $f(x)$ in einem Intervall $(\alpha,\beta)$
eingeschlossen, das keine zweite Wurzel enthält.}

\gesperrt{In dem Intervall $(\alpha,\beta)$ ist $f'(x)$ von Null verschieden und
$f''(x)$ auch von Null verschieden.}

Was die letztere Voraussetzung betrifft, dass $f''(x)$ im Intervall von
Null verschieden ist, so ist sie nur gemacht, um einfacher
auszudrückende Bedingungen für die Anwendung der Methode zu erhalten.
An sich ist ihre Brauchbarkeit bei genügender Einengung des Intervalles
davon nicht abhängig. Wenn aber $f(x)$ und $f''(x)$ keinen gemeinsamen
Theiler haben, also nicht zugleich verschwinden, so kann man die
\gesperrt{Fourier}'sche Voraussetzung durch Einengung des Intervalles
immer erfüllen, und wenn $f(x)$ und $f''(x)$ einen gemeinsamen Theiler
haben, so kann man diesen zuvor absondern und dann auf die einzelnen
Factoren von $f(x)$ die Näherungsmethoden anwenden.

Am einfachsten übersieht man die Verhältnisse in der Geometrie, wenn
man $y=f(x)$ als Gleichung einer ebenen Curve in einem rechtwinkligen
Coordinatensystem deutet.

Die Gleichung
\weq{3}{ y = f(\alpha) + (x-\alpha)f'(\alpha) } % [sic] equation number (3) repeats the earlier (3) used on the previous printed page for f(alpha)+hf'(alpha)=0; reproduced as printed, Weber restarts the numbering for this geometric sub-passage
ist die Gleichung der Curventangente in dem Punkte $\alpha, f(\alpha)$
%CONT>>
